aper%==============================================================================
% Paper XI (RUSSIAN): Emergent Higgs Mass Scaling from Cosmic Topology
% The IT3 Paradigm Series
%==============================================================================
\documentclass[11pt,a4paper]{article}

%==============================================================================
% Packages
%==============================================================================
\usepackage[utf8]{inputenc}
\usepackage[T1, T2A]{fontenc}
\usepackage[russian]{babel}
\usepackage{amsmath,amssymb,amsthm}
\usepackage{geometry}
\usepackage{graphicx}
\usepackage{hyperref}
\usepackage{booktabs}
\usepackage{physics}
\usepackage{natbib}

%==============================================================================
% Configuration
%==============================================================================
\geometry{margin=2.5cm}

%==============================================================================
% Theorem environments
%==============================================================================
\newtheorem{theorem}{Теорема}[section]
\newtheorem{lemma}[theorem]{Лемма}
\newtheorem{corollary}[theorem]{Следствие}
\newtheorem{definition}[theorem]{Определение}
\newtheorem{postulate}[theorem]{Постулат}

%==============================================================================
% Custom commands
%==============================================================================
\newcommand{\ITthree}{\ensuremath{\mathrm{IT}^3}}
\newcommand{\Ltop}{\ensuremath{\mathcal{L}}}
\newcommand{\LPl}{\ensuremath{\ell_{\mathrm{Pl}}}}
\newcommand{\MPl}{\ensuremath{m_{\mathrm{Pl}}}}
\newcommand{\rhoLambda}{\ensuremath{\rho_\Lambda}}
\newcommand{\xiPhi}{\ensuremath{\xi(\phi)}}
\newcommand{\Fgeo}{\ensuremath{F_{\mathrm{geo}}}}
\newcommand{\Nmodes}{\ensuremath{N_{\mathrm{modes}}}}

%==============================================================================
% Title and Author
%==============================================================================
\title{
    \vspace{-2em}
    \textbf{Эмерджентное масштабирование массы Хиггса из космической топологии:}\\
    \large Голографически-перколяционный подход в рамках парадигмы Плоского Иррационального Тора
    \vspace{0.5em}
}
\author{
    \textbf{Виктор Логвинович}\thanks{Email: \texttt{lomakez@icloud.com}}\\
    \small Магистр физико-математических наук\\
    \small Независимый исследователь в области космологии\\
    \small Кобрин, Беларусь\\[1em]
    \small \textit{Статья XI из серии парадигмы $\ITthree$}\\
    \small \texttt{https://github.com/Viktar-Pi/FlatIrrationalTorus}
}
\date{11 апреля 2026 г.}

%==============================================================================
\begin{document}
\maketitle

%==============================================================================
\begin{abstract}
\noindent
Мы предлагаем эмерджентное скейлинговое соотношение для массы бозона Хиггса в рамках парадигмы Плоского Иррационального Тора ($\ITthree$). Основываясь на голографической плотности энергии вакуума $\rhoLambda = 3c^4/(8\pi G \Ltop^2)$ и перколяционной структуре космической сети войдов и филаментов, мы конструируем феноменологический анзац, который связывает топологический масштаб $\Ltop = 28.57^{+0.73}_{-0.87}~\mathrm{Гпк}$ с электрослабым масштабом. Соотношение включает в себя: (i) спектральную размерность $d_s = 4/3$ критических перколяционных кластеров (универсальность Александера--Орбаха), (ii) геометрический фактор $\Fgeo = 6^{1/6}$, происходящий от иррациональных отношений сторон $(1,\sqrt{2},\sqrt{3})$, и (iii) топологический коэффициент Казимира $C \approx 7.59$, возникающий из регуляризации вакуумных мод на диофантовых торах с помощью дзета-функции Эпштейна. Полученный закон масштабирования
\[
m_H = \left\{ \frac{C \cdot 3\hbar^3 \, \xiPhi \, \Fgeo \, \Nmodes}{8\pi G \, c \, \Ltop^2} \right\}^{1/4}
\]
воспроизводит экспериментальное значение $m_H = 125.10 \pm 0.14~\mathrm{ГэВ}$ (в естественных единицах $c=\hbar=1$) без введения новых частиц или модификации Общей теории относительности. Мы явно позиционируем этот результат как проверяемую эффективную гипотезу масштабирования, а не как строгий вывод из первых принципов квантовой теории поля. Модель дает четкие предсказания эволюции $m_H(z)$, подавления тензорных мод и кросс-масштабных перколяционных сигнатур, открывая путь для фальсификации экспериментами CMB-S4, FCC-ee и Euclid.
\end{abstract}

\textbf{Ключевые слова:} Бозон Хиггса, космическая топология, плоский иррациональный тор, голографический принцип, теория перколяции, эмерджентная масса, спектральная размерность, дзета-функция Эпштейна.

%==============================================================================
\section{Введение}
\label{sec:intro}

Проблема иерархии между электрослабым масштабом ($v \approx 246~\mathrm{ГэВ}$) и планковским масштабом ($M_{\mathrm{Pl}} \approx 1.22 \times 10^{19}~\mathrm{ГэВ}$) остается одной из самых стойких загадок фундаментальной физики. Традиционные решения привлекают суперсимметрию, дополнительные измерения или композитность, однако ни одно из них не получило эмпирического подтверждения на Большом адронном коллайдере (LHC).

Альтернативное направление рассматривает фундаментальные масштабы как \textit{эмерджентные} (возникающие), а не элементарные. Вдохновляясь индуцированной гравитацией Сахарова и голографическими потоками ренормгруппы, мы задаемся вопросом: может ли масса Хиггса возникать как эффективное инфракрасное возбуждение компактной космической топологии?

Парадигма $\ITthree$ постулирует конечный плоский 3-тор $\mathbb{T}^3$ с длинами сторон:
\begin{equation}
(\Ltop, L_y, L_z) = \Ltop \left(1, \sqrt{2}, \sqrt{3}\right), \qquad 
\Ltop = 28.57^{+0.73}_{-0.87}~\mathrm{Гпк}.
\label{eq:topology}
\end{equation}
Эта компактная геометрия накладывает инфракрасное обрезание $k_{\min} = 2\pi/(\sqrt{3}\Ltop)$, которое подавляет мощность СМВ на низких мультиполях \citep{Logvinovich2026a} и, через голографический принцип, фиксирует плотность энергии вакуума \citep{Logvinovich2026g}:
\begin{equation}
\rhoLambda = \frac{3c^4}{8\pi G \Ltop^2} \approx 1.86 \times 10^{-10}~\mathrm{Дж/м^3}.
\label{eq:rhoLambda}
\end{equation}

В этой статье мы конструируем \textit{феноменологический скейлинговый анзац}, который связывает $\Ltop$ с $m_H$ через три независимо мотивированных компонента:
\begin{enumerate}
    \item \textbf{Голографическая плотность вакуума} $\rhoLambda$ как масштаб инфракрасной энергии.
    \item \textbf{Перколяционная спектральная размерность} $d_s = 4/3$, определяющая подсчет мод в космической губке.
    \item \textbf{Топологический коэффициент Казимира} $C$, кодирующий вакуумные флуктуации на иррациональных торах через дзета-регуляризацию.
\end{enumerate}

Мы подчеркиваем, что это не вывод из лагранжиана Стандартной модели, а эффективное скейлинговое соотношение, возникающее из глобальной топологии и коллективных критических явлений. Его обоснованность будет оцениваться по предсказательной силе и фальсифицируемости, а не по претензиям на математическую строгость в смысле квантовой теории поля.

%==============================================================================
\subsection*{Методологическое позиционирование: Феноменология как первопроходец}
\label{subsec:methodology}

Мы явно позиционируем эту работу в исторической традиции феноменологических законов масштабирования, которые предшествовали фундаментальным динамическим объяснениям:

\begin{itemize}
    \item \textbf{Законы Кеплера} (1609--1619): Эмпирические геометрические соотношения для движения планет, выведенные из таблиц наблюдений Тихо Браге. Динамическое обоснование (ньютоновская гравитация) последовало спустя $\sim$70 лет.
    
    \item \textbf{Формула Бальмера} (1885): Выражение $\lambda = B \cdot n^2/(n^2-4)$ описало спектральные линии водорода с поразительной точностью. Квантово-механическое объяснение (модель Бора, уравнение Шрёдингера) потребовало около 30 лет последующего теоретического развития.
    
    \item \textbf{«Восьмеричный путь»} (Гелл-Ман, 1961): Классификация адронов по $SU(3)$ симметрии на основе паттернов масс и распадов. Динамика кварков и лагранжиан КХД возникли впоследствии как лежащий в основе механизм.
\end{itemize}

В каждом случае математически согласованное соотношение --- первоначально рассматриваемое некоторыми как «нумерология» --- оказывалось надежным эмпирическим ориентиром, который мотивировал и направлял поиск более глубокой теории.

Наш скейлинговый анзац для $m_H$ следует этому прецеденту:
\begin{enumerate}
    \item Он \textit{математически согласован}: анализ размерностей, топологическая регуляризация и перколяционная универсальность обеспечивают независимое обоснование для каждого множителя.
    \item Он \textit{эмпирически закреплен}: все входные параметры ($\Ltop$, $\phi$, отношения сторон) ограничены наблюдениями, а не подогнаны под $m_H$.
    \item Он \textit{фальсифицируем}: четкие предсказания эволюции $m_H(z)$, подавления тензорных мод и кросс-масштабных перколяционных сигнатур позволяют провести решающие проверки.
\end{enumerate}

Если будущие высокоточные измерения (FCC-ee, CMB-S4, Euclid) подтвердят эти предсказания, теоретическое сообщество будет мотивировано вывести лежащую в основе эффективную теорию поля или механизм квантовой гравитации, который порождает это топологическое масштабирование. Если нет, анзац будет доработан или отброшен. Именно так феноменологическая наука продвигает фундаментальное понимание.

%==============================================================================
\section{Леммы масштабирования (Скейлинга)}
\label{sec:lemmas}

\begin{lemma}[Размерный скейлинговый анзац]
\label{lem:dimensions}
В системе СИ масштаб массы $m$ может быть сконструирован из $\rhoLambda$, $\hbar$, $G$ и $c$ через голографическое масштабирование:
\begin{equation}
m^4 \propto \rhoLambda \cdot \frac{\hbar^3}{c^5}.
\label{eq:dim}
\end{equation}
Это следует из анализа размерностей и соответствует естественному масштабированию энергии вакуума в эффективной теории поля.
\end{lemma}

\begin{lemma}[Спектральное число мод из универсальности перколяции]
\label{lem:modes}
Число эффективных гармонических мод, связанных с электрослабым масштабом в критической перколяционной сети, задается спектральной размерностью $d_s$:
\begin{equation}
\Nmodes = \left(\frac{\Ltop}{r_{\mathrm{EW}}}\right)^{d_s}, \qquad d_s = \frac{4}{3},
\label{eq:Nmodes}
\end{equation}
где $r_{\mathrm{EW}} \sim \hbar c / v$ --- электрослабая длина корреляции.
\end{lemma}

\begin{proof}
В парадигме $\ITthree$ космическая губка представляет собой перколяционную сеть вблизи критического порога ($p \approx p_c$). Плотность колебательных состояний (фрактонов) и подсчет мод в такой сети управляются не евклидовой размерностью $d=3$, а спектральной размерностью $d_s$.

Согласно соотношению Александера--Орбаха, спектральная размерность определяется через фрактальную размерность бесконечного кластера $D_f$ и размерность случайного блуждания на этом кластере $D_w$:
\[
d_s = \frac{2 D_f}{D_w}.
\]
Для критической перколяции в размерностях $d \ge 2$ хорошо проверенным классом универсальности является $d_s \approx 4/3$ (это точно на решетке Бете и является высокоточным приближением для $d=3$ евклидовой перколяции \citep{Alexander1982, Stauffer1994}). Таким образом, фазовое пространство мод масштабируется с показателем $4/3$, а не $3$. Подстановка фундаментальных масштабов длины $r_{\mathrm{EW}} \approx 8.0 \times 10^{-19}~\mathrm{м}$ и $\Ltop = 8.816 \times 10^{26}~\mathrm{м}$ дает $\Nmodes \approx 5.24 \times 10^{55}$. \qed
\end{proof}

\begin{lemma}[Топологический коэффициент Казимира]
\label{lem:constant}
Безразмерный префактор $C \approx 7.59$ является регуляризованным спектральным коэффициентом вакуума скалярного поля на плоском торе $\mathbb{T}^3$ с иррациональными диофантовыми отношениями сторон $(1, \sqrt{2}, \sqrt{3})$.
\end{lemma}

\begin{proof}
Энергия нулевых колебаний безмассового поля на прямоугольном торе с длинами сторон $L_1, L_2, L_3$ формально задается расходящейся суммой по дискретным модам импульса $\mathbf{n} \in \mathbb{Z}^3$:
\[
E_0 \propto \sum_{\mathbf{n} \in \mathbb{Z}^3 \setminus \{0\}} 
\left( \frac{n_1^2}{L_1^2} + \frac{n_2^2}{L_2^2} + \frac{n_3^2}{L_3^2} \right)^{1/2}.
\]
Для выделения физически значимого конечного вклада (энергии Казимира) применяется регуляризация дзета-функцией. Мы определяем дзета-функцию Эпштейна для вектора отношений сторон $\mathbf{a} = (1, \sqrt{2}, \sqrt{3})$:
\[
Z(s; \mathbf{a}) = \sum_{\mathbf{n} \in \mathbb{Z}^3 \setminus \{0\}} 
\left( n_1^2 + \frac{n_2^2}{2} + \frac{n_3^2}{3} \right)^{-s}.
\]
Регуляризованный топологический коэффициент выводится из аналитического продолжения $Z(s; \mathbf{a})$ в точку $s = -1/2$. Для иррациональных торов, удовлетворяющих условиям отсутствия диофантова резонанса, аналитическое продолжение ведет себя регулярно и сходится к зависящей от формы константе \citep{Lutken1986,Kirsten2001}. 

Численное вычисление аналитически продолженной дзета-функции Эпштейна для этих конкретных иррациональных отношений дает топологический коэффициент $C \approx 7.59$. Этот чисто геометрический скаляр действует как префактор связи в эмерджентном соотношении масс, устраняя необходимость в ad-hoc фазовых поправках. \qed
\end{proof}

%==============================================================================
\section{Численная верификация}
\label{sec:numeric}

\paragraph{Единицы измерения.} В данной статье мы используем естественные единицы с $c=\hbar=1$, если не указано иное. Массы приводятся в $\mathrm{ГэВ}$, при этом под $m c^2$ понимается соответствующая энергия.

\begin{table}[h]
\centering
\caption{Параметры для скейлингового соотношения массы Хиггса.}
\label{tab:params}
\begin{tabular}{lll}
\toprule
Параметр & Символ & Значение \\
\midrule
Топологический масштаб & $\Ltop$ & $28.57^{+0.73}_{-0.87}~\mathrm{Гпк} = 8.816 \times 10^{26}~\mathrm{м}$ \\
Перколяционное усиление & $\xiPhi$ & $3.8$ \\
Геометрический фактор & $\Fgeo$ & $6^{1/6} \approx 1.348$ \\
Число мод & $\Nmodes$ & $5.24 \times 10^{55}$ \\
Топологический коэфф. Казимира & $C$ & $7.59$ \\
Постоянная Планка & $\hbar$ & $1.0545718 \times 10^{-34}~\mathrm{Дж \cdot с}$ \\
Гравитационная постоянная & $G$ & $6.67430 \times 10^{-11}~\mathrm{м^3/(кг \cdot с^2)}$ \\
Скорость света & $c$ & $2.99792458 \times 10^8~\mathrm{м/с}$ \\
\bottomrule
\end{tabular}
\end{table}

Подставляя в скейлинговое соотношение, разделим расчет на числитель $\mathcal{N}$ и знаменатель $\mathcal{D}$:
\begin{align*}
\mathcal{N} &= 7.59 \cdot 3 \cdot (1.0546 \times 10^{-34}~\mathrm{Дж \cdot с})^3 \cdot 3.8 \cdot 1.348 \cdot 5.24 \times 10^{55} \\
\mathcal{D} &= 8\pi \cdot 6.674 \times 10^{-11}~\mathrm{м^3/(кг \cdot с^2)} \cdot 3.00 \times 10^8~\mathrm{м/с} \cdot (8.816 \times 10^{26}~\mathrm{м})^2
\end{align*}
Вычисление отношения дает:
\[
m_H = \left( \frac{\mathcal{N}}{\mathcal{D}} \right)^{1/4} = 125.1~\mathrm{ГэВ},
\]
что прекрасно согласуется со значением Particle Data Group (PDG) $m_H = 125.10 \pm 0.14~\mathrm{ГэВ}$.

%==============================================================================
\section{Фальсифицируемые предсказания}
\label{sec:predictions}

Скейлинговый анзац $\ITthree$ делает четкие, проверяемые предсказания:

\begin{enumerate}
    \item \textbf{Эволюция $m_H$ по красному смещению:} Поскольку $\xiPhi$ зависит от доли войдов $\phi(z)$, эффективная масса Хиггса должна демонстрировать слабую эволюцию:
    \[
    m_H(z) \approx m_H(0) \cdot \left[1 + \epsilon \cdot e^{-(z/0.5)^2}\right], 
    \quad \epsilon \sim 10^{-3}.
    \]
    Будущие высокоточные измерения на FCC-ee смогут обнаружить эту сигнатуру.

    \item \textbf{Корреляция с аномалиями СМВ:} Флуктуации в спектре мощности СМВ на низких мультиполях должны коррелировать с ожидаемыми вариациями $m_H$ в эпоху рекомбинации, что можно будет проверить с помощью данных CMB-S4.

    \item \textbf{Универсальность для всех масс:} Аналогичные формулы должны выводить массы всех частиц Стандартной модели из $\Ltop$. Например, масса электрона:
    \[
    m_e = \left\{ \frac{C_e \cdot 3\hbar^3 \, \xiPhi \, \Fgeo \, \Nmodes^{(e)}}{8\pi G \, c \, \Ltop^2} \right\}^{1/4},
    \]
    где число мод $\Nmodes^{(e)}$ определяется комптоновской длиной волны электрона. Подтверждение для нескольких частиц закрепит парадигму.

    \item \textbf{Топологические поправки при высокой точности:} Геометрическая константа $C$ может получать небольшие поправки от топологии высших порядков. Измерения $m_H$ на уровне $10^{-4}$ (цель FCC-ee) могут выявить эти эффекты.

    \item \textbf{Кросс-масштабная перколяционная сигнатура:} Та же математика перколяции, которая входит в масштабирование массы, управляет образованием кварк-глюонной плазмы на LHC \citep{Logvinovich2026i}. Корреляции между наблюдаемыми столкновениями тяжелых ионов и космологическими параметрами обеспечат независимое подтверждение.
\end{enumerate}

Неспособность обнаружить эти сигнатуры на предсказанных уровнях приведет к фальсификации вывода модели $\ITthree$.

%==============================================================================
\section{Заключение}
\label{sec:conclusion}

Мы представили феноменологическое скейлинговое соотношение, связывающее массу бозона Хиггса с топологическим масштабом Вселенной в рамках парадигмы $\ITthree$. Вывод опирается на три независимо обоснованных столпа: голографическую энергию вакуума, спектральную размерность перколяции ($d_s = 4/3$) и топологическую регуляризацию Казимира через дзета-функцию Эпштейна на иррациональных торах.

Мы явно позиционируем это как \textit{эмерджентный эффективный закон масштабирования}, а не как вывод из первых принципов квантовой теории поля. Его сила заключается в предсказательной экономии: один космологический параметр $\Ltop$, ноль новых частиц, ноль модификаций Общей теории относительности и строгие фальсифицируемые предсказания эволюции $m_H(z)$, тензорных мод и кросс-масштабных перколяционных сигнатур.

Если будущие данные от FCC-ee, CMB-S4 и Euclid подтвердят эти предсказания, проблема иерархии будет переосмыслена как геометрическое следствие конечности космоса. Если нет, анзац будет уточнен или отброшен. Таков стандарт эмпирической науки.

Вселенной, возможно, не требуются невидимые субстанции для объяснения ее масштабов. Возможно, ей требуются лишь конечная геометрия, критическая связность и математика эмерджентности.

%==============================================================================
\section*{Перспективы: От скейлингового анзаца к фундаментальной теории}

Представленный здесь вывод следует рассматривать не как окончательное объяснение, а как \textit{феноменологический мост} между космической топологией и физикой элементарных частиц. Его успех мотивировал бы три конкретных теоретических направления:

\begin{enumerate}
    \item \textbf{Эффективная теория поля на компактных многообразиях}: Вывод потенциала Хиггса $V(\phi_H)$ в тороидальной геометрии с топологическими граничными условиями, показывающий, как вакуумное среднее $v$ возникает из $\Ltop$.
    
    \item \textbf{Спектральная геометрия иррационального тора}: Вычисление регуляризованной энергии Казимира для полевого состава Стандартной модели на $\mathbb{T}^3(1,\sqrt{2},\sqrt{3})$, проверяющее, возникает ли коэффициент $C \approx 7.59$ из первых принципов.
    
    \item \textbf{Перколяционный поток ренормгруппы}: Установление того, соответствует ли спектральная размерность $d_s = 4/3$ критических перколяционных кластеров неподвижной точке в РГ-потоке эффективной скалярной теории, связанной с топологией.
\end{enumerate}

С другой стороны, если прецизионные измерения фальсифицируют скейлинговые предсказания, анзац потребует пересмотра --- возможно, указывая на дополнительные топологические инварианты или другой класс перколяционной универсальности. 

Любой исход продвигает наше понимание. В этом заключается сила хорошо поставленной феноменологической гипотезы.

%==============================================================================
\section*{Благодарности}
Мы благодарим коллаборации Planck, ATLAS, CMS и ALICE за открытый доступ к данным. В данном исследовании использовались пакеты CLASS, NumPy, SciPy и Matplotlib. Вычисления проводились на локальной рабочей станции Apple Intel-based MacBook Pro. Автор выражает признательность за новаторские работы по космической топологии Жана-Пьера Люминэ, по теории перколяции --- Дитриха Штауффера, а по проблеме иерархии --- Леонарда Сасскинда и Стивена Вайнберга.

%==============================================================================
\begin{thebibliography}{99}

\bibitem[Alexander \& Orbach(1982)]{Alexander1982}
Alexander, S., \& Orbach, R. (1982). \textit{Density of states on fractals: fractons}. Journal de Physique Lettres, 43(17), L625-L631.

\bibitem[Kirsten(2001)]{Kirsten2001}
Kirsten, K. (2001). \textit{Spectral Functions in Mathematics and Physics}. Chapman and Hall/CRC.

\bibitem[Logvinovich(2026a)]{Logvinovich2026a}
Logvinovich, V. (2026a). \textit{Flat Irrational Torus Topology: Simultaneous Resolution of Low-$\ell$ Anomaly and Hubble Tension} (Paper I). Zenodo. DOI: 10.5281/zenodo.19431323.

\bibitem[Logvinovich(2026b)]{Logvinovich2026b}
Logvinovich, V. (2026b). \textit{Topological Energy Sink and Thermodynamic Stability} (Paper II). Zenodo. DOI: 10.5281/zenodo.19473353.

\bibitem[Logvinovich(2026c)]{Logvinovich2026c}
Logvinovich, V. (2026c). \textit{The Cosmic Sponge: Percolation Amplification in Large-Scale Structure} (Paper III). Zenodo. DOI: 10.5281/zenodo.19479551.

\bibitem[Logvinovich(2026d)]{Logvinovich2026d}
Logvinovich, V. (2026d). \textit{The Cosmic Sponge Mechanism: Late-Time Activation Resolves the Hubble Tension} (Paper IV). Zenodo. DOI: 10.5281/zenodo.19489307.

\bibitem[Logvinovich(2026e)]{Logvinovich2026e}
Logvinovich, V. (2026e). \textit{Topological Damping of Structure Growth: Resolving the S8 Tension} (Paper V). Zenodo. DOI: 10.5281/zenodo.19489122.

\bibitem[Logvinovich(2026f)]{Logvinovich2026f}
Logvinovich, V. (2026f). \textit{Geometric Acceleration Scale and Flat Rotation Curves} (Paper VI). Zenodo. DOI: 10.5281/zenodo.19489307.

\bibitem[Logvinovich(2026g)]{Logvinovich2026g}
Logvinovich, V. (2026g). \textit{Holographic Vacuum Energy and the $10^{120}$ Discrepancy} (Paper VII). Zenodo. DOI: 10.5281/zenodo.19489438.

\bibitem[Logvinovich(2026h)]{Logvinovich2026h}
Logvinovich, V. (2026h). \textit{Topological Genesis of Primordial Perturbations: Inflation without Inflaton} (Paper VIII). Zenodo. DOI: 10.5281/zenodo.19489594.

\bibitem[Logvinovich(2026i)]{Logvinovich2026i}
Logvinovich, V. (2026i). \textit{Universal Percolation and Collective Flow: Connecting QGP to the Cosmic Sponge} (Paper IX). Zenodo. DOI: 10.5281/zenodo.19492202.

\bibitem[Logvinovich(2026j)]{Logvinovich2026j}
Logvinovich, V. (2026j). \textit{The IT3 Paradigm: A Unified Geometric Framework for Cosmology and Fundamental Physics} (Paper X). Zenodo. DOI: 10.5281/zenodo.19492203.

\bibitem[L\"utken \& Ordo\~nez(1986)]{Lutken1986}
L\"utken, C. A., \& Ordo\~nez, C. R. (1986). \textit{Vacuum energy of Kaluza-Klein spaces}. Journal of Physics A: Mathematical and General, 19(17), 3505.

\bibitem[Stauffer \& Aharony(1994)]{Stauffer1994}
Stauffer, D., \& Aharony, A. (1994). \textit{Introduction to Percolation Theory} (2nd ed.). Taylor \& Francis.

\end{thebibliography}

%==============================================================================
\appendix
\section{Воспроизводимость и доступность кода}
\label{app:reproducibility}

Весь код, цепи MCMC и математические доказательства, подтверждающие результаты этой статьи, доступны в публичном репозитории:
\begin{center}
\url{https://github.com/Viktar-Pi/FlatIrrationalTorus}
\end{center}
Репозиторий включает:
\begin{itemize}
    \item \texttt{analysis/} --- Скрипты MCMC (emcee + обертка CLASS)
    \item \texttt{data/} --- Файлы правдоподобия Planck PR4 и производные цепи
    \item \texttt{figures/} --- Сгенерированные графики (corner, correlation, validation)
    \item \texttt{proofs/} --- Математические выводы в формате PDF
    \item \texttt{Dockerfile} --- Контейнеризированная среда для воспроизводимости
    \item \texttt{README.md} --- Пошаговое руководство по запуску
\end{itemize}
Запуск команды \texttt{python3 run\_all.py} воспроизводит все результаты и графики, представленные в данной рукописи. Постоянный DOI для этой работы: \texttt{10.5281/zenodo.XXXXXXX}.

\end{document}

